\subsection{P-20. Two-pair evaluation for translated consecutive linear blocks}
\label{P-20}

\paragraph{Statement.}
Fix distinct nodes $a_1,\ldots,a_n$ over any field $E$.
Let $f\ne0$ have ROABP width at most three in some variable
order. Suppose its variables have a partition into nonempty
consecutive blocks $B_1,\ldots,B_d$ in that order such that
$f$ is homogeneous of degree one in each entire block.
Assume also the formal common-translation identity
\[
                  f(X+t\mathbf1)=f(X).
 \tag{P20.1}
\]
Then balanced rational two-pair evaluation hits $f$. This
holds for arbitrary $d$, every characteristic, and unknown
order and blocks. For total degree $d$, the finite hitting
set has at most
\[
                 (n+1)^2\binom{d+3}{3}
 \tag{P20.2}
\]
points over a working field of size at least
$\max\{2n+1,d+1\}$, with I-40's extension-input convention.

For $d=4$, the translation hypothesis can be omitted:
every nonzero width-three polynomial homogeneous of degree
one in each of four consecutive blocks is hit by the same
map, with at most $35(n+1)^2$ points. The cubic case is
already covered without any block promise by P-19.

\paragraph{Inputs and scope.}
Write $\mathcal I_k$ for the prime ideal of balanced rational
degree-$k$ evaluation at the fixed nodes. We use the already
proved absence of linear equations in $\mathcal I_1$, the
inclusion $\mathcal I_2\subseteq\mathcal I_1$, and pole
extraction $\partial_i\mathcal I_2\subseteq\mathcal I_1$.
The pole-extraction proof in I-37 applies to every multilinear
degree; it is not restricted to cubics. We use exact tensor
rank factorization at the block cuts (P-B17, Oseledets,
Theorem~2.1, pp.~2298--2299), rather than assume that the
original univariate factors already have block width two.
All block ranks and within-block ranks below are intrinsic
coefficient ranks of the scalar polynomial.

\paragraph{1. Separate translations force block-cut rank at most two.}
The coefficient of $t$ in (P20.1) is
\[
                   \sum_{b=1}^d\sum_{i\in B_b}\partial_i f=0.
\]
The $d$ summands have different block multidegrees, so each
is zero. Since the polynomial has degree one in each entire
block, this gives the full separate common-translation
identity in each block, in every characteristic.

Remove unused variables. Each block still has essential
variables, and in fact at least two: a single essential
variable in a degree-one block would contradict its separate
translation identity. Fix a block $B_b$ and write
\[
             f=\sum_{i\in B_b}X_i H_i,
             \qquad \sum_{i\in B_b}H_i=0.
 \tag{P20.3}
\]
Let $i_*$ be its last essential position. At the cut after
the whole block, let the coefficient rank be $r_b$.
The coefficient rows supplied by the earlier $H_i$ already
span that entire row space, because $H_{i_*}$ is their
negative sum. At the cut just before $i_*$, these earlier
terms form a degree-$b$ left block of rank $r_b$.
The nonzero term $X_{i_*}H_{i_*}$ forms a separate
degree-$(b-1)$ left block of positive rank. The two blocks
have disjoint row and column degrees. Hence
\[
                         r_b+1\le3,
                         \qquad r_b\le2.
 \tag{P20.4}
\]
Only the internal block cuts will be needed below.

\paragraph{2. A smallest counterexample has two-dimensional block cores.}
Suppose a counterexample to the first assertion exists.
Extend scalars to an algebraic closure, which preserves
formal ideal membership and all coefficient ranks. Choose
one with the fewest blocks, retaining consecutive order and
the stated translation property. A block-cut rank of one
would factor $f=gh$ into nonconstant polynomials on the
prefix and suffix blocks. Both factors are separately
homogeneous of degree one in their retained blocks, and
inherit width at most three by specializing the other
factor to a nonzero value. They also inherit common
translation invariance: the identity for the product makes
the ratio for each factor a scalar in the translation
parameter, and its highest homogeneous part makes that
scalar one. Primality of $\mathcal I_2$ puts one factor
in $\mathcal I_2$, contradicting the minimal block count.
Thus every internal block-cut rank is exactly two.

The cases of one or two blocks have total degree at most
two and cannot be counterexamples by the established
minimum degree of $\mathcal I_2$. We may assume $d\ge3$.
Exact rank factorization of the block coefficient tensor
now gives
\[
 f=L(X_{B_1})^{\mathsf T}
       M_2(X_{B_2})\cdots M_{d-1}(X_{B_{d-1}})
       R(X_{B_d}),
 \tag{P20.5}
\]
where $L,R$ have two independent linear-form entries and
each $M_b$ is a $2\times2$ matrix of homogeneous linear
forms in its own block. The prefix and suffix coefficient
functions at every internal block cut are independent.
In particular, separate block translation implies
\[
 M_b=\sum_{i\in B_b}X_i M_{b,i},\qquad
                       \sum_{i\in B_b}M_{b,i}=0.
 \tag{P20.6}
\]
This follows from independence on the two disjoint sides,
so the intervening constant matrix is uniquely determined.

Every interior $M_b$ has generic rank two. Otherwise its
nonzero homogeneous linear entries have a polynomial
rank-one factorization with one constant vector factor:
the two nonnegative degrees in a primitive factorization
sum to one. A constant column factor would make the cut
before the block rank one; a constant row factor would
do the same after the block. Both contradict minimality.
Equivalently, a linear space of matrices of rank at most
one has either a common image line or a common row line;
the same conclusion follows directly from its rank-one
coefficient matrices.

\paragraph{3. The within-block width bound factors every determinant.}
Fix an interior block and suppress its first subscript.
Order its essential variables as $X_1,\ldots,X_h$, with
coefficient matrices $M_1,\ldots,M_h$. Every matrix is
nonzero, and $\sum_iM_i=0$.
For a proper prefix $\{1,\ldots,k\}$ of this block, put
\[
 \begin{aligned}
 \rho_k&=\dim\sum_{i\le k}\operatorname{row}(M_i),\\
 \kappa_k&=\dim\sum_{i>k}\operatorname{col}(M_i).
 \end{aligned}
\]
Independence of the two preceding coefficient functions
and the two following ones makes the within-block
coefficient rank exactly
\[
                         \rho_k+\kappa_k\le3.
 \tag{P20.7}
\]
Indeed the prefix-variable terms have one larger left
degree and rank equal to the stacked-row rank $\rho_k$;
the suffix-variable terms have rank equal to the
concatenated-column rank $\kappa_k$. These are disjoint
homogeneous degree blocks.

Let $p$ be the first index at which the accumulated row
space has dimension two. It exists because the block
matrix has generic rank two. One has $1<p<h$. If $p=h$,
the zero sum would put $M_h$ in the common earlier row
line, a contradiction. If $p=1$, then (P20.7) makes all
remaining columns lie in one line; their sum would make
$M_1$ have rank at most one, again a contradiction.

The matrices before $p$ have one common row line, and
those after $p$ have one common column line. Choose
nonzero constant vectors $a,b$ and write
\[
 M_i=u_i a^{\mathsf T}\quad(i<p),\qquad
 M_i=b v_i^{\mathsf T}\quad(i>p).
\]
The zero sum fixes the remaining coefficient matrix.
Consequently
\[
 \begin{aligned}
 M(X)&=U(X)a^{\mathsf T}+bV(X)^{\mathsf T},\\
 U(X)&=\sum_{i<p}(X_i-X_p)u_i,\qquad
 V(X)=\sum_{i>p}(X_i-X_p)v_i,\\
 \det M(X)&=\det[U(X),b]\,
                  \det\begin{pmatrix}a^{\mathsf T}\\V(X)^{\mathsf T}\end{pmatrix}.
 \end{aligned}
 \tag{P20.8}
\]
The last two factors are homogeneous linear forms.
Both are nonzero because $M$ has generic rank two.
This is a consequence of the actual within-block cut
ranks and the translation identity, not an assumption
about simultaneous diagonalization or individual factors.

\paragraph{4. A one-pole center preserves the block pencils and first derivatives suffice.}
Use the common one-pole center
\[
                  \alpha_i=t+\frac{z}{a_i-y}
\]
over its rational function field $K=\overline E(t,z,y)$.
No nonzero linear form vanishes at this center, by the
absence of linear equations in $\mathcal I_1$.
Equation (P20.8) therefore makes every interior block
matrix invertible at $\alpha$. The two endpoint vectors
in (P20.5) remain nonzero as well.

Translate $X=\alpha+Y$ and use constant cut gauges over
$K$ to normalize the interior block constants to the
identity. The endpoints are affine-linear functions of
their block variables, with nonzero constants $u_0,v_0$.
They can be absorbed into affine-linear matrices of
constant term $I_2$. For example, if $(u_0)_j\ne0$,
\[
 u(Y)=u_0\left(I_2+
       \frac{e_j}{(u_0)_j}(u(Y)-u_0)\right),
\]
and the column analogue uses any nonzero coordinate of
$v_0$. Thus the whole translated polynomial has the form
\[
 f(\alpha+Y)=u_0
   \prod_{b=1}^{d}\left(I_2+
                  \sum_{i\in B_b}Y_i N_i\right)v_0.
 \tag{P20.9}
\]
The matrices $N_i$ are over $K$; the variables in distinct
blocks remain disjoint.

Since $f\in\mathcal I_2$, both $f(\alpha)=0$ and
$\partial_i f(\alpha)=0$ for every $i$. In (P20.9) these
identities say
\[
                  u_0v_0=0,\qquad u_0N_i v_0=0.
 \tag{P20.10}
\]
The annihilator of the nonzero column $v_0$ in the
two-dimensional row space is a line. It contains the
nonzero row $u_0$, so (P20.10) gives
$u_0N_i\in K u_0$ for every $i$. Iterating this equality
through any ordered product of coefficient matrices
shows that every coefficient of (P20.9) is zero. This
contradicts $f\ne0$ and proves the first assertion.
The denominator-safe selector and amplitude interpolation
of I-40 gives (P20.2).

\paragraph{5. The four-block assertion does not need a translation promise.}
Suppose a homogeneous multilinear quartic of the stated
four-block type were a width-three annihilator of balanced
two-pair evaluation. Every positive-power coefficient of
$f(X+t\mathbf1)$ in $t$ belongs to $\mathcal I_2$ and
has degree at most three. If one were nonzero, the highest
such coefficient would retain every coefficient-cut rank
upper bound by the leading-parameter specialization used
in I-40. P-19 excludes its cubic case, and the minimum
degree of $\mathcal I_2$ excludes smaller degrees.
All those coefficients therefore vanish, proving (P20.1).
The first assertion now gives a contradiction. Substituting
$d=4$ in (P20.2) yields the stated $35(n+1)^2$ bound.
\hfill$\square$

\paragraph{Remaining scope.}
The consecutive blocks must be linear as entire blocks;
arbitrary multiaffine blocks and arbitrary linear changes
of coordinates are not covered. For more than four blocks
the common-translation identity remains an explicit
hypothesis. The determinant-factorization proof is used
only after intrinsic block ranks have become two, and
does not assert preservation of general rank-two reset
words or full matrix-image cones. The construction gives
a restricted polynomial sampler, not a polynomial bound
in growing ROABP width. No historical priority is claimed.
